\begin{ex}
Trong các số: $8, 3, 10, 6, 9$, số lớn nhất là số nào?
\shortans{10}
\loigiai{
So sánh các số đã cho: $3 < 6 < 8 < 9 < 10$.\nSố lớn nhất trong các số trên là số $10$.
\textbf{Đáp số:} $10$.
}
\end{ex}

\begin{ex}
Trong các số: $7, 2, 5, 0, 9$, số bé nhất là số nào?
\shortans{0}
\loigiai{
Số bé nhất trong dãy số tự nhiên từ 0 đến 10 là số $0$.
\textbf{Đáp số:} $0$.
}
\end{ex}

\begin{ex}
Tìm số tự nhiên $x$ thỏa mãn: $6 < x < 8$.
\shortans{7}
\loigiai{
Số tự nhiên nằm giữa 6 và 8 là số liền sau của 6 và liền trước của 8, đó là số $7$.
\textbf{Đáp số:} $7$.
}
\end{ex}

\begin{ex}
Có bao nhiêu số tự nhiên $x$ thỏa mãn: $4 < x < 9$?
\shortans{4}
\loigiai{
Các số tự nhiên lớn hơn 4 và bé hơn 9 là: $5, 6, 7, 8$.\nĐếm các số trên có tất cả $4$ số.
\textbf{Đáp số:} $4$.
}
\end{ex}

\begin{ex}
Cho các số: $9, 4, 7, 2, 10$. Sắp xếp các số theo thứ tự từ bé đến lớn, số đứng ở vị trí thứ ba là số nào?
\shortans{7}
\loigiai{
Sắp xếp các số từ bé đến lớn: $2 < 4 < 7 < 9 < 10$.\nSố đứng ở vị trí thứ ba là số $7$.
\textbf{Đáp số:} $7$.
}
\end{ex}

\begin{ex}
Cho các số: $10, 5, 8, 3, 7$. Sắp xếp các số theo thứ tự từ lớn đến bé, số đứng ở vị trí thứ hai là số nào?
\shortans{8}
\loigiai{
Sắp xếp các số từ lớn đến bé: $10 > 8 > 7 > 5 > 3$.\nSố đứng ở vị trí thứ hai là số $8$.
\textbf{Đáp số:} $8$.
}
\end{ex}

\begin{ex}
Tìm số lớn nhất có một chữ số bé hơn 9.
\shortans{8}
\loigiai{
Các số có một chữ số bé hơn 9 là: $0, 1, 2, 3, 4, 5, 6, 7, 8$.\nSố lớn nhất trong các số này là số $8$.
\textbf{Đáp số:} $8$.
}
\end{ex}

\begin{ex}
Tìm số bé nhất có một chữ số lớn hơn 5.
\shortans{6}
\loigiai{
Các số có một chữ số lớn hơn 5 là: $6, 7, 8, 9$.\nSố bé nhất trong các số này là số $6$.
\textbf{Đáp số:} $6$.
}
\end{ex}

\begin{ex}
Lấy số lớn nhất trong các số $6, 2, 10, 7$ trừ đi số bé nhất trong các số đó, ta được kết quả bằng bao nhiêu?
\shortans{8}
\loigiai{
Trong các số đã cho:\n- Số lớn nhất là số $10$.\n- Số bé nhất là số $2$.\nHiệu của số lớn nhất và số bé nhất là: $10 - 2 = 8$.
\textbf{Đáp số:} $8$.
}
\end{ex}

\begin{ex}
Có bao nhiêu số tự nhiên lớn hơn 2 và bé hơn hoặc bằng 8?
\shortans{6}
\loigiai{
Các số tự nhiên đó là: $3, 4, 5, 6, 7, 8$.\nĐếm các số trên có tất cả $6$ số.
\textbf{Đáp số:} $6$.
}
\end{ex}

\begin{ex}
Tìm số tự nhiên $a$ sao cho: $a + 1 = 9$. Giá trị của $a$ là bao nhiêu?
\shortans{8}
\loigiai{
Muốn tìm số hạng $a$, ta lấy tổng trừ đi 1: $a = 9 - 1 = 8$.
\textbf{Đáp số:} $8$.
}
\end{ex}

\begin{ex}
Tìm số tự nhiên $b$ sao cho: $10 - b = 3$. Giá trị của $b$ là bao nhiêu?
\shortans{7}
\loigiai{
Muốn tìm số trừ $b$, ta lấy số bị trừ là 10 trừ đi hiệu là 3: $b = 10 - 3 = 7$.
\textbf{Đáp số:} $7$.
}
\end{ex}

\begin{ex}
Điền số thích hợp vào ô trống để biểu thức đúng: $5 < \dots < 7$.
\shortans{6}
\loigiai{
Số tự nhiên đứng giữa 5 và 7 là số $6$.
\textbf{Đáp số:} $6$.
}
\end{ex}

\begin{ex}
Cho ba số $a, b, c$. Biết $a > b$ và $b > c$. Nếu $a = 9$ và $c = 7$, hỏi $b$ bằng bao nhiêu?
\shortans{8}
\loigiai{
Theo đề bài: $7 < b < 9$.\nSố tự nhiên nằm giữa 7 và 9 là số $8$. Vậy $b = 8$.
\textbf{Đáp số:} $8$.
}
\end{ex}

\begin{ex}
Tìm số tự nhiên $x$ lớn nhất thỏa mãn: $x < 10$.
\shortans{9}
\loigiai{
Các số tự nhiên bé hơn 10 là: $0, 1, 2, \dots, 9$.\nSố lớn nhất trong các số này là số $9$.
\textbf{Đáp số:} $9$.
}
\end{ex}

\begin{ex}
Tìm số tự nhiên $y$ bé nhất thỏa mãn: $y > 6$.
\shortans{7}
\loigiai{
Các số tự nhiên lớn hơn 6 là: $7, 8, 9, 10, \dots$\nSố bé nhất trong các số này là số $7$.
\textbf{Đáp số:} $7$.
}
\end{ex}

\begin{ex}
Có bao nhiêu số tự nhiên có một chữ số lớn hơn 3?
\shortans{6}
\loigiai{
Các số có một chữ số lớn hơn 3 là: $4, 5, 6, 7, 8, 9$.\nĐếm các số này có tất cả $6$ số.
\textbf{Đáp số:} $6$.
}
\end{ex}

\begin{ex}
Tổng của số lớn nhất có một chữ số và số bé nhất có một chữ số là bao nhiêu?
\shortans{9}
\loigiai{
Số lớn nhất có một chữ số là số $9$.\nSố bé nhất có một chữ số là số $0$.\nTổng của hai số là: $9 + 0 = 9$.
\textbf{Đáp số:} $9$.
}
\end{ex}

\begin{ex}
Hiệu giữa số bé nhất có hai chữ số và số lớn nhất có một chữ số bằng bao nhiêu?
\shortans{1}
\loigiai{
Số bé nhất có hai chữ số là số $10$.\nSố lớn nhất có một chữ số là số $9$.\nHiệu của hai số là: $10 - 9 = 1$.
\textbf{Đáp số:} $1$.
}
\end{ex}

\begin{ex}
Cho dãy số: $2, 4, 6, 8, 10$. Số đứng ở vị trí thứ tư trong dãy là số nào?
\shortans{8}
\loigiai{
Đếm theo thứ tự từ trái sang phải:\n- Thứ nhất là 2; thứ hai là 4; thứ ba là 6; thứ tư là 8; thứ năm là 10.\nVậy số đứng ở vị trí thứ tư là số $8$.
\textbf{Đáp số:} $8$.
}
\end{ex}

\begin{ex}
So sánh hai phép tính: $4 + 5$ và $3 + 6$. Kết quả của mỗi phép tính đều bằng bao nhiêu?
\shortans{9}
\loigiai{
Ta có: $4 + 5 = 9$ và $3 + 6 = 9$.\nKết quả của hai phép tính đều bằng $9$.
\textbf{Đáp số:} $9$.
}
\end{ex}

\begin{ex}
Số nào trong các số $6, 7, 8, 9$ khi cộng với 2 thì được kết quả lớn hơn 9?
\shortans{8}
\loigiai{
Thử từng số:\n- $6 + 2 = 8 < 9$.\n- $7 + 2 = 9$ (bằng 9).\n- $8 + 2 = 10 > 9$ (chọn).\nSố cần tìm là số $8$.
\textbf{Đáp số:} $8$.
}
\end{ex}

\begin{ex}
Tìm số tự nhiên $m$ biết rằng: $8 - m > 5$ và $m$ là số tự nhiên lớn nhất có thể. Giá trị của $m$ là bao nhiêu?
\shortans{2}
\loigiai{
Ta có: $8 - 1 = 7 > 5$; $8 - 2 = 6 > 5$; $8 - 3 = 5$ (không lớn hơn 5).\nVậy các số $m$ thỏa mãn là 0, 1, 2. Giá trị lớn nhất của $m$ là số $2$.
\textbf{Đáp số:} $2$.
}
\end{ex}

\begin{ex}
Trong một cuộc thi chạy, bạn Nam về đích sau bạn An nhưng trước bạn Bình. Biết bạn An về đích thứ nhất, bạn Bình về đích thứ 3. Hỏi bạn Nam về đích thứ mấy?
\shortans{2}
\loigiai{
Bạn Nam về đích sau người thứ nhất và trước người thứ ba, vậy bạn Nam về đích ở vị trí thứ $2$.
\textbf{Đáp số:} $2$.
}
\end{ex}

\begin{ex}
Hiệu của số đứng liền sau số 8 và số đứng liền trước số 4 là bao nhiêu?
\shortans{6}
\loigiai{
Số đứng liền sau số 8 là số $9$.\nSố đứng liền trước số 4 là số $3$.\nHiệu của hai số là: $9 - 3 = 6$.
\textbf{Đáp số:} $6$.
}
\end{ex}
